\chapter*{Towns and Hamlets}\stepcounter{chapter}\phantomsection\addcontentsline{toc}{chapter}{Towns and Hamlets}

\section*{Battle in the Skies}\phantomsection\addcontentsline{toc}{section}{Battle in the Skies}
{\entryfont Supernatural Hazard \& Combat\\
\textbf{Recommended Level:} 5-7\\
\textbf{Difficulty:} Hard (Deadly with Complications)}

\subsection*{Setup}
{\entryfont In the outskirts of a small settlement near Dundee, the aftermath of the Unicorn Invasion of Dundee lingers in a far more insidious form. The souls of those slain have not passed on - they swirl together into a dreadful \hyperref[monster:Sluagh]{\LinkFont{Sluagh}} that haunts the skies above the town.

\textbf{If the party arrives during the day}, the town appears eerily still. Smoke-blackened ruins and abandoned homes dominate the landscape, but a handful of survivors remain hidden. These desperate individuals emerge cautiously, pleading for aid. They speak of a "storm of screaming dead" that descends with the coming dusk, taking anyone caught outside.}
\begin{DndReadAloud}
	The village lies in ruin - charred beams, collapsed roofs, and an oppressive silence. From behind broken walls, gaunt faces peer out. One survivor stumbles forward, voice trembling: "You can't stay out when the sun sets... it comes for us. Please... help us before night falls."
\end{DndReadAloud}
{\noindent\entryfont \textbf{If the party arrives at dusk or night}, the haunting is already under way. The sky churns with a mass of spectral forms, and the air is filled with agonized wailing. The Sluagh actively hunts, descending without warning to seize victims and drag them into the sky. The party enters directly into chaos, with little time for preparation.}
\begin{DndReadAloud}
	A chorus of distant wails drifts across the ruined rooftops, growing into a deafening cacophony. Above, a writhing mass of dark, translucent figures churns like a living storm. Suddenly, a lone figure is yanked screaming into the sky, their body vanishing into the mass. The air grows cold, and the ground itself seems to recoil - as if the dead refuse to rest.
\end{DndReadAloud}

\subsection*{Mechanics}
\subsubsection*{Day-Time Preparation Phase}
{\entryfont During daylight, the Sluagh remains dormant or dispersed, its presence only faintly perceptible (whispers, cold spots, shadows). This grants the party time to prepare. Skill checks such as Religion or Arcana can reveal the nature of the entity and how to weaken it. Survivors can assist in gathering materials for rituals, setting traps, or securing buildings.}
\subsubsection*{Night-Time Assault Phase}
{\entryfont At dusk, two Sluagh fully manifest and begin their assault. They hover 30-60 ft. above ground, descending unpredictably to grapple targets before lifting them skyward. Strong winds impose disadvantage on long-range attacks, and the Sluagh's incorporeal movement allows it to ignore most terrain and cover.}

\subsection*{Possible Solutions}
\subsubsection*{Direct Confrontation}
{\entryfont Engaging the Sluagh requires careful coordination to prevent allies from being isolated and carried off. Forcing it into sunlight or grounding it (via control effects such as restraining magic) significantly reduces its threat. Focused damage while it is bloodied may disrupt its cohesion, weakening its attacks.}
\subsubsection*{Ritualistic Release}
{\entryfont If the party chooses to perform a ritual to disperse or appease the Sluagh, treat it as a multi-stage skill challenge conducted either during the day (preparation) or under pressure at night. The ritual requires 5 total successes before 3 failures. Each attempt represents roughly 1 round (during combat) or 10-15 minutes (during the day).\\
\textbf{Primary Checks:}
\begin{itemize}
	\item Religion (DC 15)
	\item Arcana (DC 15)
	\item Performance (DC 14)
\end{itemize}
\textbf{Secondary / Support Checks:}
\begin{itemize}
	\item Medicine (DC 13)
	\item History (DC 13)
	\item Survival (DC 13)
\end{itemize}
Upon reaching 5 successes, the ritual completes and the Sluagh will fully be dispersed.}

\subsection*{Possible Complications}
\subsubsection*{Reinforcements from Undead Horde}
{\entryfont The commotion may draw remnants of the nearby undead forces - specifically \hyperref[monster:UndeadSoldier]{\LinkFont{Undead Soldier}} patrols or even a roaming \hyperref[monster:CorruptedUnicorn]{\LinkFont{Corrupted Unicorn}}. These arrive mid-encounter, splitting the party's attention and increasing pressure dramatically.}
\subsubsection*{Civilian Panic}
{\entryfont Terrified townsfolk may run into danger, forcing the party to choose between offence and rescue, potentially exposing themselves to aerial grabs.}

\subsection*{Outcomes}
\subsubsection*{Sluagh Dispersed (Ritual Success)}
{\entryfont The swirling mass dissolves into faint lights that drift skyward, leaving behind an eerie calm. The town is freed, and grateful survivors offer what little they have. The area may become sanctified ground.}
\subsubsection*{Sluagh Destroyed (Combat Victory)}
{\entryfont The entity collapses violently, releasing a shock wave of necrotic energy before fading. Fragments of spiritual essence crystallize briefly before vanishing.}
\subsubsection*{Rewards}
{\entryfont\begin{itemize}
	\item A well-maintained weapon (e.g. common enchantment)
	\item A charm or brooch worth 15 GP
	\item A signet ring that could grant advantage on social checks
\end{itemize}}